% ============================================================
%  chapter03.tex — 第三章 平面奇点指数
%  对应原书第三章。
%  引入向量场旋转数与奇点指数的拓扑不变量，并给出指数计算的
%  代数公式（Cauchy 指标、Bendixson 公式等）。
% ============================================================

\newcommand{\chThreeFigurePlaceholder}[1]{%
  \begin{center}
    \fbox{\parbox{0.72\textwidth}{\centering\vspace{2.7cm}图~#1\\（图片占位）\vspace{2.7cm}}}
  \end{center}%
}

\setchapterimage{}
\setchapterintro{}
\chapter{平面奇点指数}

% ============================================================
%  §1  连续向量场的旋转数
% ============================================================
\section{连续向量场的旋转数}

给定平面 $\mathbb{R}^2$ 上的连续向量场
\[
  A(x,y)=\bigl(X(x,y),Y(x,y)\bigr),\qquad (x,y)\in\mathbb{R}^2.
\]
设 $\mathscr{L}\subset\mathbb{R}^2$ 是一条有向的闭分段光滑曲线，且
$A(x,y)$ 在 $\mathscr{L}$ 上非退化，即
\[
  X^2(x,y)+Y^2(x,y)\ne0,\qquad (x,y)\in\mathscr{L}.
\]
作单位化映射
\[
  T(x,y)=\dfrac{A(x,y)}{\lVert A(x,y)\rVert}
  =\left(
    \dfrac{X(x,y)}{\sqrt{X^2(x,y)+Y^2(x,y)}},
    \dfrac{Y(x,y)}{\sqrt{X^2(x,y)+Y^2(x,y)}}
  \right).
\]
它将 $\mathscr{L}$ 映入单位圆周。映射 $T$ 的度 $\mu(T,\mathscr{L})$
称为向量场 $A$ 沿 $\mathscr{L}$ 的\emph{旋转数}，记为
\[
  \gamma(A,\mathscr{L})=\mu(T,\mathscr{L}). \tag{1.1}
\]

若 $A(x,y)$ 光滑，则
\[
  \gamma(A,\mathscr{L})
  =\dfrac{1}{2\pi}\oint_{\mathscr{L}}
  \mathrm{d}\arctan\dfrac{Y}{X}. \tag{1.2}
\]

设 $D\subset\mathbb{R}^2$ 是有界闭区域，其边界为
\[
  \partial D=\bigcup_{i=1}^n\partial D_i,
\]
并取 $\partial D_i$ 的正向，使区域在行进方向的左侧。定义
\[
  \gamma(A,\partial D)=\sum_{i=1}^n\gamma(A,\partial D_i).
\]

\chThreeFigurePlaceholder{3.1}
\chThreeFigurePlaceholder{3.2}
\chThreeFigurePlaceholder{3.3}

\noindent\textbf{性质~1}\quad
设 $D_1,D_2$ 是两个内部没有公共点的有界闭区域，且 $A$ 在
$\partial D_1\cup\partial D_2$ 上非退化，则
\[
  \gamma(A,\partial(D_1\cup D_2))
  =\gamma(A,\partial D_1)+\gamma(A,\partial D_2). \tag{1.3}
\]

\noindent\textbf{性质~2}\quad
若 $A$ 在有界闭区域 $D$ 上连续且非退化，则
\[
  \gamma(A,\partial D)=0.
\]
事实上，当 $D$ 单连通时，上式由连续向量场的光滑逼近即得。一般地，
设 $\partial D=\bigcup_{i=1}^n\partial D_i$，将 $D$ 分成
$n-1$ 个单连通区域 $G_i$，如图~3.4 所示。由单连通情形和性质~1，
便有
\[
  \gamma(A,\partial D)=\sum_{i=1}^{n-1}\gamma(A,\partial G_i)=0.
\]

\chThreeFigurePlaceholder{3.4}

\begin{definition}{}{homotopy-vector-field}
设 $D\subset\mathbb{R}^2$ 是有界闭区域，$A_0(x,y),A_1(x,y)$
是 $\partial D$ 上的连续、非退化向量场。若存在依赖于参数
$\lambda$ 的连续向量场，也可看成连续映射
\[
  A(x,y,\lambda):\partial D\times[0,1]\longrightarrow\mathbb{R}^2,
\]
且
\[
  A(x,y,0)=A_0(x,y),\qquad A(x,y,1)=A_1(x,y),
\]
并且当 $(x,y)\in\partial D$、$\lambda\in[0,1]$ 时
$A(x,y,\lambda)$ 恒不为零，则称 $A(x,y,\lambda)$ 为从
$A_0(x,y)$ 到 $A_1(x,y)$ 的连续形变；此时称 $A_0(x,y)$ 与
$A_1(x,y)$ 同伦。
\end{definition}

\begin{theorem}{}{homotopy-transitivity}
若 $A_0(x,y),A_1(x,y)$ 都与 $A_*(x,y)$ 同伦，则
$A_0(x,y)$ 与 $A_1(x,y)$ 也同伦。
\end{theorem}

\begin{proof}
设从 $A_i(x,y)$ 到 $A_*(x,y)$ 的连续、非退化形变分别为
$\overline{A}_i(x,y,\lambda)$，$i=0,1$。令
\[
  A(x,y,\lambda)=
  \begin{cases}
    \overline{A}_0(x,y,2\lambda),&0\leq\lambda\leq\dfrac12,\\
    \overline{A}_1(x,y,2-2\lambda),&\dfrac12<\lambda\leq1,
  \end{cases}
\]
则 $A(x,y,\lambda)$ 即为从 $A_0(x,y)$ 到 $A_1(x,y)$ 的连续、
非退化形变。
\end{proof}

\noindent\textbf{性质~3}\quad
设 $D\subset\mathbb{R}^2$ 是有界闭区域，则 $\partial D$ 上的同伦
向量场有相同的旋转数。事实上，$\gamma(A(x,y,\lambda),\partial D)$
一方面为整数，另一方面关于 $\lambda$ 连续，故只能为常数。

\begin{theorem}{}{same-rotation-nonopposite}
设 $D\subset\mathbb{R}^2$ 是有界闭区域，在 $\partial D$ 上有两个
非退化、连续向量场 $A_0(x,y)$ 和 $A_1(x,y)$。若对一切
$(x,y)\in\partial D$，$A_0(x,y)$ 和 $A_1(x,y)$ 的方向恒不相反，即
\[
  \dfrac{A_0(x,y)}{\lVert A_0(x,y)\rVert}
  \ne-\dfrac{A_1(x,y)}{\lVert A_1(x,y)\rVert},
  \qquad (x,y)\in\partial D,
\]
则
\[
  \gamma(A_0,\partial D)=\gamma(A_1,\partial D). \tag{1.5}
\]
\end{theorem}

\begin{proof}
构造连续形变
\[
  A(x,y,\lambda)=(1-\lambda)A_0(x,y)+\lambda A_1(x,y),
  \qquad (x,y)\in\partial D,\quad 0\leq\lambda\leq1.
\]
若它在某点退化，则
\[
  (1-\lambda_0)A_0(x_0,y_0)=-\lambda_0A_1(x_0,y_0),
\]
从而 $A_0(x_0,y_0)$ 和 $A_1(x_0,y_0)$ 的方向相反，矛盾。因此该形变
非退化，结论由性质~3 得证。
\end{proof}

\begin{definition}{}{continuous-perturbation}
设 $D\subset\mathbb{R}^2$ 是有界闭区域，设 $A_0(x,y)$ 是
$\partial D$ 上的连续、非退化向量场。向量场
$\Delta A_0(x,y)$ 称作 $A_0(x,y)$ 的连续微小扰动，若
$\Delta A_0(x,y)$ 在 $\partial D$ 上连续，且
\[
  \lVert\Delta A_0(x,y)\rVert<\lVert A_0(x,y)\rVert,
  \qquad (x,y)\in\partial D.
\]
\end{definition}

\begin{theorem}{}{small-perturbation-rotation}
设 $D\subset\mathbb{R}^2$ 是有界闭区域。若 $A_0(x,y)$ 是
$\partial D$ 上的连续、非退化向量场，$\Delta A_0(x,y)$ 是
$A_0(x,y)$ 的连续微小扰动，令
\[
  A_1(x,y)=A_0(x,y)+\Delta A_0(x,y),
\]
则 $A_1(x,y)$ 在 $\partial D$ 上也非退化，且
\[
  \gamma(A_0,\partial D)=\gamma(A_1,\partial D). \tag{1.6}
\]
\end{theorem}

\begin{proof}
构造从 $A_0(x,y)$ 到 $A_1(x,y)$ 的连续形变
\[
  A(x,y,\lambda)=A_0(x,y)+\lambda\Delta A_0(x,y),
  \qquad (x,y)\in\partial D,\quad0\leq\lambda\leq1.
\]
由微小扰动的定义，
\[
  \lVert A(x,y,\lambda)\rVert
  \geq\lVert A_0(x,y)\rVert-\lambda\lVert\Delta A_0(x,y)\rVert
  >(1-\lambda)\lVert\Delta A_0(x,y)\rVert\geq0.
\]
故 $A_0(x,y)$ 与 $A_1(x,y)$ 同伦，由性质~3 即得式~(1.6)。
\end{proof}

% ============================================================
%  §2  平面奇点指数
% ============================================================
\section{平面奇点指数}

对于由连续向量场
\[
  A(x,y)=\bigl(X(x,y),Y(x,y)\bigr)
\]
所确定的平面微分方程
\[
  \dfrac{\mathrm{d}x}{\mathrm{d}t}=X(x,y),\qquad
  \dfrac{\mathrm{d}y}{\mathrm{d}t}=Y(x,y), \tag{2.1}
\]
如果 $M_0(x_0,y_0)$ 为方程~(2.1) 的孤立奇点
\[
  A(x_0,y_0)=0,
\]
则存在充分小的正数 $R$，使圆盘 $S_R=S(M_0,R)$ 内部无异于
$M_0$ 的奇点。对任意 $0<r_1<r_2<R$，环域
$S_{r_2}\setminus S_{r_1}$ 内无奇点，因而
\[
  \gamma(A,\partial S_{r_2})
  =\gamma\bigl(A,\partial(S_{r_2}\setminus S_{r_1})\bigr)
  +\gamma(A,\partial S_{r_1})
  =\gamma(A,\partial S_{r_1}).
\]
可见，对任意 $0<r<R$，$\gamma(A,\partial S_r)$ 是不依赖于
$r$ 的整数。

\begin{definition}{}{index-of-singularity}
考虑微分方程~(2.1)，设 $X(x,y),Y(x,y)\in C^0$，又设
$M_0(x_0,y_0)$ 是其孤立奇点。对一切 $0<r<R$，
$\gamma(A,\partial S_r)$ 是常数，称它为方程~(2.1) 的奇点
$M_0(x_0,y_0)$ 的\emph{指数}，记作
\[
  J_{M_0}(A)\quad\text{或}\quad J_{M_0}(X,Y).
\]
当方程~(2.1) 只有唯一奇点时，简记为 $J(A)$ 或 $J(X,Y)$。
\end{definition}

连续向量场的奇点有下列几个性质。

\noindent\textbf{性质~1}\quad
若有界单连通或多连通区域 $D$ 的内部仅包含方程~(2.1) 的有限个奇点
$M_1,M_2,\ldots,M_k$，且在 $\partial D$ 上无奇点，则
\[
  \gamma(A,\partial D)=\sum_{i=1}^kJ_{M_i}(A). \tag{2.2}
\]

\begin{proof}
取 $r_i>0$ 充分小，使
\[
  S_{r_i}=S(M_i,r_i)\subset D,\qquad
  S_{r_i}\cap S_{r_j}=\varnothing\quad(i\ne j).
\]
由第一节的性质~1、2，
\[
  \gamma(A,\partial D)
  =\gamma\left(A,\partial\left(D\setminus\bigcup_{i=1}^kS_{r_i}\right)\right)
   +\sum_{i=1}^k\gamma(A,\partial S_{r_i})
  =\sum_{i=1}^kJ_{M_i}(A).
\]
\end{proof}

\noindent\textbf{性质~2}\quad
若方程~(2.1) 只有孤立奇点，则方程的任何闭轨线 $\mathscr L$ 所围成的
区域 $D$ 内部所有奇点的指数和为 $1$。

\noindent\textbf{性质~3}\quad
方程~(2.1) 的闭轨线 $\mathscr L$ 所围成的区域 $D$ 内部一定包含奇点。

\noindent\textbf{性质~4}\quad
\[
  J_{M_0}(aX,Y)=\operatorname{sgn}a\cdot J_{M_0}(X,Y),
  \qquad a\ne0.
\]

\noindent\textbf{性质~5}\quad
\[
  J_{M_0}(X,Y)=-J_{M_0}(Y,X).
\]

\noindent\textbf{性质~6}\quad
\[
  J_{M_0}(X+Y,Y)=J_{M_0}(X,Y).
\]

为了说明性质~4、5、6 在计算奇点指数时的用处，下面应用公式~(1.2)
及这些性质来求线性奇点的指数。虽然下面的例题~1--5 都是习题~4 的特例，
仍将它们单独加以讨论。

\begin{example}{}{index-node}
方程
\[
  \dfrac{\mathrm{d}x}{\mathrm{d}t}=x,\qquad
  \dfrac{\mathrm{d}y}{\mathrm{d}t}=y
\]
的奇点 $O(0,0)$ 是临界结点，求它的指数。

\begin{solution}
由公式~(1.2)，取 $\mathscr L:x^2+y^2=1$，则
\[
  J(x,y)=\dfrac{1}{2\pi}\oint_{\mathscr L}
  \mathrm{d}\tan^{-1}\dfrac{y}{x}
  =\dfrac{1}{2\pi}\int_0^{2\pi}\mathrm{d}\theta=1.
\]
\end{solution}
\end{example}

\begin{example}{}{index-center}
方程
\[
  \dfrac{\mathrm{d}x}{\mathrm{d}t}=-y,\qquad
  \dfrac{\mathrm{d}y}{\mathrm{d}t}=x
\]
的奇点 $O(0,0)$ 是中心，求它的指数。

\begin{solution}
由例~2.1 及奇点指数的性质~4、5，
\[
  J(-y,x)=-J(y,x)=J(x,y)=1.
\]
\end{solution}
\end{example}

\begin{example}{}{index-saddle-node}
求方程
\[
  \dfrac{\mathrm{d}x}{\mathrm{d}t}=\lambda_1x,\qquad
  \dfrac{\mathrm{d}y}{\mathrm{d}t}=\lambda_2y,\qquad
  \lambda_1\lambda_2\ne0
\]
的奇点 $O(0,0)$ 的指数。

\begin{solution}
由例~2.1 及性质~4、5，
\[
  J(\lambda_1x,\lambda_2y)
  =\operatorname{sgn}(\lambda_1\lambda_2)J(x,y)
  =\operatorname{sgn}(\lambda_1\lambda_2).
\]
当 $\lambda_1\lambda_2<0$ 时，$O(0,0)$ 为鞍点，指数为 $-1$；
当 $\lambda_1\lambda_2>0$ 时，$O(0,0)$ 为结点，指数为 $+1$。
\end{solution}
\end{example}

\begin{example}{}{index-focus}
方程
\[
  \dfrac{\mathrm{d}x}{\mathrm{d}t}=\mu_1x-\mu_2y,\qquad
  \dfrac{\mathrm{d}y}{\mathrm{d}t}=\mu_2x+\mu_1y,\qquad
  \mu_1\mu_2\ne0
\]
的奇点 $O(0,0)$ 是焦点，求它的指数。

\begin{solution}
由例~2.1 及性质~4、5、6，
\begin{align*}
  J(\mu_1x-\mu_2y,\mu_2x+\mu_1y)
  &=\operatorname{sgn}(\mu_1\mu_2)
    J(\mu_1^2x-\mu_1\mu_2y,\mu_2^2x+\mu_1\mu_2y)\\
  &=J(y,-x)=J(x,y)=1.
\end{align*}
\end{solution}
\end{example}

\begin{example}{}{index-degenerate-node}
方程
\[
  \dfrac{\mathrm{d}x}{\mathrm{d}t}=\lambda_1x,\qquad
  \dfrac{\mathrm{d}y}{\mathrm{d}t}=\xi x+\lambda_1y,\qquad
  \xi\lambda_1\ne0
\]
的奇点 $O(0,0)$ 是退化结点，求它的指数。

\begin{solution}
由例~2.1 及性质~4、5、6，
\[
  J(\lambda_1x,\xi x+\lambda_1y)=J(y,-x)=J(x,y)=1.
\]
\end{solution}
\end{example}

如果 $X(x,y)$ 和 $Y(x,y)$ 是 $x,y$ 的解析函数，则称
$A(x,y)$ 为解析向量场。对于由解析向量场所确定的方程组~(2.1) 的奇点
$M_0$，可将 $X,Y$ 在 $M_0$ 上作 Taylor 展开：
\[
  X(x,y)=P_m(x-x_0,y-y_0)+\widetilde X(x-x_0,y-y_0),
\]
\[
  Y(x,y)=Q_n(x-x_0,y-y_0)+\widetilde Y(x-x_0,y-y_0).
\]
这里 $m$ 次齐次多项式 $P_m$ 和 $n$ 次齐次多项式 $Q_n$ 分别是
$X(x,y)$ 和 $Y(x,y)$ 的展开式中的最低次项之和；
$\widetilde X,\widetilde Y$ 分别是所有次数高于 $m,n$ 的各项之和。此时
\[
  \dfrac{\mathrm{d}x}{\mathrm{d}t}=P_m(x,y),\qquad
  \dfrac{\mathrm{d}y}{\mathrm{d}t}=Q_n(x,y) \tag{2.3}
\]
称作方程组~(2.1) 在奇点 $M_0=(x_0,y_0)$ 上的主方程。方程~(2.3) 的
奇点 $x=y=0$ 对应于方程组~(2.1) 的奇点 $M_0$。通常 $m=n$ 时称为
齐次奇点；当 $m$ 和 $n$ 不都为 $1$ 时，方程组~(2.1) 的奇点 $M_0$
称作高次奇点。

\begin{theorem}{}{principal-system-index}
设方程组~(2.1) 中 $X(x,y),Y(x,y)$ 在奇点 $M_0=(x_0,y_0)$ 上解析。
若主方程~(2.3) 的奇点 $x=y=0$ 孤立，那么方程组~(2.1) 的奇点 $M_0$
也孤立，并且
\[
  J_{M_0}(X,Y)=J(P_m,Q_n).
\]
\end{theorem}

\begin{proof}
构造从向量场
\[
  \bigl(P_m(x-x_0,y-y_0),Q_n(x-x_0,y-y_0)\bigr)
\]
到 $(X(x,y),Y(x,y))$ 的连续形变
\[
  A_\lambda(x,y)=\bigl(X_\lambda(x,y),Y_\lambda(x,y)\bigr),
\]
其中
\[
  X_\lambda=P_m(x-x_0,y-y_0)+\lambda\widetilde X(x-x_0,y-y_0),
\]
\[
  Y_\lambda=Q_n(x-x_0,y-y_0)+\lambda\widetilde Y(x-x_0,y-y_0),
  \qquad 0\leq\lambda\leq1.
\]
若结论不成立，则存在异于 $M_0$ 且以 $M_0$ 为极限的点列
$(x_k,y_k)$ 和数列 $\lambda_k\in[0,1]$，使
$A_{\lambda_k}(x_k,y_k)=0$。令
\[
  x_k-x_0=r_k\cos\theta_k,\qquad
  y_k-y_0=r_k\sin\theta_k,
\]
则 $r_k\to0$。取收敛子列，令
\[
  (\cos\theta_k,\sin\theta_k)\to
  (\cos\bar\theta,\sin\bar\theta),\qquad
  \lambda_k\to\bar\lambda.
\]
将等式除以适当的 $r_k^m,r_k^n$ 后令 $k\to\infty$，得到
\[
  P_m(\cos\bar\theta,\sin\bar\theta)
  =Q_n(\cos\bar\theta,\sin\bar\theta)=0.
\]
于是直线
\[
  \dfrac{x}{\cos\bar\theta}=\dfrac{y}{\sin\bar\theta}
\]
上的所有点都是主方程~(2.3) 的奇点，这与其原点为孤立奇点的假设矛盾。
故对充分小的圆 $S(M_0,R)$，$A_\lambda$ 在其边界上恒非退化。由同伦不变性，
即得
\[
  J_{M_0}(X,Y)=J(P_m,Q_n).
\]
\end{proof}

% ============================================================
%  §3  Cauchy指标
% ============================================================
\section{Cauchy指标}

给定两个 $t$ 的实系数多项式
\[
  u(t)=\sum_{i=0}^k a_it^{k-i},\qquad
  v(t)=\sum_{i=0}^l b_it^{l-i}. \tag{3.1}
\]
设分式函数 $v(t)/u(t)$ 的所有变号间断点为
\[
  t_1,t_2,\ldots,t_r,t'_1,t'_2,\ldots,t'_{r'},
\]
其中
\[
  \lim_{t\to t_i-0}\dfrac{v(t)}{u(t)}=-\infty,\qquad
  \lim_{t\to t_i+0}\dfrac{v(t)}{u(t)}=+\infty
  \quad(i=1,2,\ldots,r),
\]
\[
  \lim_{t\to t'_i-0}\dfrac{v(t)}{u(t)}=+\infty,\qquad
  \lim_{t\to t'_i+0}\dfrac{v(t)}{u(t)}=-\infty
  \quad(i=1,2,\ldots,r').
\]
则称整数 $r-r'$ 为分式函数 $v(t)/u(t)$ 的 Cauchy 指标，表示为
\[
  N\bigl(u(t),v(t)\bigr)\quad\text{或}\quad N(u,v).
\]

\begin{example}{}{cauchy-index-example}
设
\[
  v(t)=t^2,\qquad
  u(t)=(t-2)^2(t-4)^3(t-5)(t-6),
\]
求 $N(u,v)$。

\begin{solution}
$v(t)/u(t)$ 的变号间断点为 $t=4,5,6$。在 $t=4,5,6$ 的左右极限
依次为
\[
  -\infty,+\infty;\qquad +\infty,-\infty;\qquad -\infty,+\infty.
\]
故
\[
  N(u,v)=2-1=1.
\]
\end{solution}
\end{example}

下面证明 Cauchy 指标 $N(u,v)$ 的几个性质。

\noindent\textbf{性质~1}\quad
如果式~(3.1) 中 $a_0\ne0$，且 $u,v$ 无公共实根，那么
$N(u,v)-k$ 为偶数，即 $N(u,v)$ 与整数 $k$ 奇偶性相同。

事实上，此时 $u$ 是 $v/u$ 的变号间断点的单根（或奇重根），
变号间断点的总个数 $r+r'$ 与 $N(u,v)=r-r'$ 奇偶性相同；而
$u(t)$ 的单根（或奇重根）总个数与其次数 $k$ 奇偶性相同。

\noindent\textbf{性质~2}\quad
设式~(3.1) 中 $a_0b_0\ne0$，则
\[
  \dfrac{1}{\pi}\int_{-\infty}^{+\infty}
  \mathrm{d}\tan^{-1}\dfrac{v(t)}{u(t)}
  =
  \begin{cases}
    \operatorname{sgn}(a_0b_0)-N(u,v),&l-k\text{ 为正奇数},\\
    -N(u,v),&l-k\text{ 不为正奇数}.
  \end{cases} \tag{3.2}
\]

\begin{proof}
设 $v(t)/u(t)$ 的所有间断点为 $\lambda_i$ $(i=1,2,\ldots,s)$，
并记 $t\to\lambda_i-0$ 和 $t\to\lambda_i+0$ 时的符号分别为
$\varepsilon_i^-$、$\varepsilon_i^+$；又以
$\varepsilon^+,\varepsilon^-$ 分别表示 $t\to+\infty,t\to-\infty$
时的符号。如图~3.5 所示，在相邻间断点之间，
\[
  \dfrac{1}{\pi}\int_{\lambda_i}^{\lambda_{i+1}}
  \mathrm{d}\tan^{-1}\dfrac{v(t)}{u(t)}
  =\dfrac12(\varepsilon_{i+1}^--\varepsilon_i^+),
  \qquad i=1,2,\ldots,s-1. \tag{3.3}
\]
当 $l\leq k$ 时，两端有有限且相等的极限。若
$\lim_{t\to\pm\infty}\tan^{-1}(v/u)=\theta$，则
\[
  \dfrac{1}{\pi}\left(
    \int_{-\infty}^{\lambda_1}+\int_{\lambda_s}^{+\infty}
  \right)\mathrm{d}\tan^{-1}\dfrac{v(t)}{u(t)}
  =\dfrac12(\varepsilon_1^--\varepsilon_s^+). \tag{3.4}
\]
若 $l>k$，则相应的两段积分为
\[
  \dfrac{1}{\pi}\int_{-\infty}^{\lambda_1}
  \mathrm{d}\tan^{-1}\dfrac{v(t)}{u(t)}
  =\dfrac12(\varepsilon_1^--\varepsilon^-),
\]
\[
  \dfrac{1}{\pi}\int_{\lambda_s}^{+\infty}
  \mathrm{d}\tan^{-1}\dfrac{v(t)}{u(t)}
  =\dfrac12(\varepsilon^+-\varepsilon_s^+). \tag{3.5}
\]
当 $l-k$ 为正偶数时 $\varepsilon^+=\varepsilon^-$；当 $l-k$
为正奇数时，
\[
  \varepsilon^+=\operatorname{sgn}(a_0b_0),\qquad
  \varepsilon^-=-\operatorname{sgn}(a_0b_0).
\]
将上述各式相加，并注意
\[
  N(u,v)=\dfrac12\sum_{i=1}^s(\varepsilon_i^--\varepsilon_i^+),
\]
即得式~(3.2)。
\end{proof}

\chThreeFigurePlaceholder{3.5}

\noindent\textbf{性质~3}\quad
设式~(3.1) 中 $a_0b_0\ne0$，则
\[
  N(u,v)+N(v,u)
  =\operatorname{sgn}\bigl[(1-(-1)^{l-k})a_0b_0\bigr].
\]

\noindent\textbf{性质~4}\quad
\[
  \lvert N(u,v)\rvert\leq\min(k,l+1).
\]
事实上，$v(t)/u(t)$ 的变号间断点为 $u(t)$ 的根，其总个数不超过
$u(t)$ 的最大次数 $k$，故 $|N(u,v)|\leq k$；再利用性质~3 即得结论。

下面推导两个计算 Cauchy 指标的常用公式。

\noindent\textbf{公式~1}\quad
如果 $u,v$ 无公共实根，且 $u(t)$ 无实重根。令
$\lambda_i$ $(i=1,2,\ldots,s)$ 是 $u(t)$ 的所有实根，那么
\[
  N(u,v)=\sum_{i=1}^s
  \operatorname{sgn}\bigl(u'(\lambda_i)v(\lambda_i)\bigr).
\]

对于式~(3.1) 中给定的多项式 $u$ 与 $v$ 辗转相除（余式改变符号），即
\[
  v=w_0u-u_1,\qquad u=w_1u_1-u_2,
\]
\[
  u_{i-1}=w_iu_i-u_{i+1}\quad(i=2,3,\ldots,r),\qquad
  u_{r+1}=0.
\]
令多项式 $u_i$ 的首项为
$\alpha_it^{m_i}$ $(\alpha_i\ne0,\ 1\leq i\leq r)$，则有

\noindent\textbf{公式~2}\quad
\[
  N(u,v)=\sum_{i=1}^r
  \dfrac{(-1)^{m_i}-1}{2}\operatorname{sgn}\alpha_i.
\]

\begin{proof}
因为
\[
  \dfrac{v(t)}{u(t)}=w_0(t)-\dfrac{u_1(t)}{u(t)},
\]
故
\[
  N(u,v)=-N(u,u_1).
\]
又 $u$ 与 $u_1$ 的次数差为 $m_1$，首项系数比为 $\alpha_1$。由性质~3，
\[
  N(u,u_1)-N(u_1,u)
  =\dfrac{(-1)^{m_1}-1}{2}\operatorname{sgn}\alpha_1.
\]
同理，
\[
  N(u_1,u)-N(u_2,u_1)
  =\dfrac{(-1)^{m_2}-1}{2}\operatorname{sgn}\alpha_2,
\]
\[
  N(u_i,u_{i-1})-N(u_{i+1},u_i)
  =\dfrac{(-1)^{m_{i+1}}-1}{2}\operatorname{sgn}\alpha_{i+1}.
\]
又 $u_{r+1}=0$，故 $N(u_r,u_{r-1})=0$。将各式相加即得公式~2。
\end{proof}

\begin{example}{}{cauchy-index-euclidean}
设
\[
  u(t)=t^3+t^2+t-1,\qquad v(t)=t^2+2t+3,
\]
求 $N(u,v)$ 及 $N(v,u)$，并验证性质~1、3、4。

\begin{solution}
辗转相除得
\[
  v=w_0u-u_1,\qquad w_0=0,
\]
\[
  u=w_1u_1-u_2,\qquad w_1=-t+1,
\]
\[
  u_1=w_2u_2,\qquad
  w_2=\dfrac12t^2+t+\dfrac32.
\]
因此
\[
  N(u,v)
  =\dfrac{(-1)^1-1}{2}\operatorname{sgn}(-1)
   +\dfrac{(-1)^2-1}{2}\operatorname{sgn}\left(\dfrac12\right)
  =1.
\]
同样可得
\[
  N(v,u)=0.
\]
\end{solution}
\end{example}

% ============================================================
%  §4  齐次方程孤立奇点指数的有理计算
% ============================================================
\section{齐次方程孤立奇点指数的有理计算}

对于型如式~(2.3) 的微分方程，其中
$P_m,Q_n$ 分别为 $x,y$ 的 $m$ 次齐次多项式和 $n$ 次齐次多项式。
当奇点 $x=y=0$ 孤立时，即 $P_m,Q_n$ 无公因子时，
$P_m(1,0),Q_n(1,0)$ 不全为零。此时关于奇点指数有下列定理。

\begin{theorem}{}{homogeneous-index}
设 $x=y=0$ 为式~(2.3) 的孤立奇点，则
\[
J(P_m,Q_n)=
\begin{cases}
0, & m+n\text{ 为奇数};\\[4pt]
N\bigl(P_m(t,1),Q_n(t,1)\bigr),
& m+n\text{ 为偶数，且 }P_m(1,0)\ne0;\\[4pt]
-N\bigl(Q_n(t,1),P_m(t,1)\bigr),
& m+n\text{ 为偶数，且 }Q_n(1,0)\ne0.
\end{cases}
\]
\end{theorem}

\begin{proof}
奇点的指数等于向量场 $(P_m,Q_n)$ 沿单位圆周旋转一周的旋转数，故
\[
\begin{aligned}
J(P_m,Q_n)
&=\dfrac{1}{2\pi}\int_{-\pi}^{\pi}
\mathrm{d}\arctan\dfrac{Q_n(\cos\theta,\sin\theta)}
{P_m(\cos\theta,\sin\theta)}\\
&=\dfrac{1}{2\pi}\left(\int_{-\pi}^{0}+\int_{0}^{\pi}\right)
\mathrm{d}\arctan\dfrac{Q_n(\cos\theta,\sin\theta)}
{P_m(\cos\theta,\sin\theta)}.
\end{aligned}
\]

当 $m+n$ 为奇数时，在第一个积分中令 $\tau=\theta+\pi$，便有
\[
\begin{aligned}
J(P_m,Q_n)
&=-\dfrac{1}{2\pi}\int_{0}^{\pi}
\mathrm{d}\arctan\dfrac{Q_n(\cos\theta,\sin\theta)}
{P_m(\cos\theta,\sin\theta)}\\
&\quad+\dfrac{1}{2\pi}\int_{0}^{\pi}
\mathrm{d}\arctan\dfrac{Q_n(\cos\theta,\sin\theta)}
{P_m(\cos\theta,\sin\theta)}=0.
\end{aligned}
\]

当 $m+n$ 为偶数时，
\[
J(P_m,Q_n)=\dfrac{1}{\pi}\int_{0}^{\pi}
\mathrm{d}\arctan\dfrac{Q_n(\cos\theta,\sin\theta)}
{P_m(\cos\theta,\sin\theta)}.
\]
令 $r=\dfrac{m-n}{2}$，并令 $t=\dfrac{\cos\theta}{\sin\theta}$，则
\[
\begin{aligned}
J(P_m,Q_n)
&=-\dfrac{1}{\pi}\int_{-\infty}^{+\infty}
\mathrm{d}\arctan\left[
\sin^{n-m}\theta\,
\dfrac{Q_n\left(\dfrac{\cos\theta}{\sin\theta},1\right)}
{P_m\left(\dfrac{\cos\theta}{\sin\theta},1\right)}
\right]\\
&=-\dfrac{1}{\pi}\int_{-\infty}^{+\infty}
\mathrm{d}\arctan\left[
(t^2+1)^r\dfrac{Q_n(t,1)}{P_m(t,1)}
\right].
\end{aligned}
\]

若 $P_m(1,0)\ne0$，则 $P_m(t,1)$ 的次数为 $m$，而
$(t^2+1)^rQ_n(t,1)$ 的次数不超过 $2r+n=m$。由第3节的式
~(3.2)，并注意到 $(t^2+1)^r$ 不改变符号，
得
\[
J(P_m,Q_n)=N\bigl(P_m(t,1),Q_n(t,1)\bigr).
\]
若 $Q_n(1,0)\ne0$，由对称性以及
$J(P_m,Q_n)=-J(Q_n,P_m)$，可得
\[
J(P_m,Q_n)=-N\bigl(Q_n(t,1),P_m(t,1)\bigr).\qedhere
\]
\end{proof}

当式~(2.1) 在奇点 $M_0$ 的主方程~(2.3) 具有孤立原点时，利用定理
\ref{thm:principal-system-index}和定理\ref{thm:homogeneous-index}，
便可求得 $M_0$ 的指数。

\begin{theorem}{}{index-bound}
当 $m+n$ 为偶数且 $x=y=0$ 为式~(2.3) 的孤立奇点时，
$J(P_m,Q_n)$ 与 $m$ 或 $n$ 同奇偶，并且
\[
\left|J(P_m,Q_n)\right|\leq\min\{m,n\}. \tag{4.1}
\]
\end{theorem}

\begin{proof}
若 $P_m(1,0)\ne0$，则 $P_m(t,1)$ 的次数为 $m$。由定理
\ref{thm:homogeneous-index}及第3节 Cauchy 指标的性质，
$J(P_m,Q_n)$ 与 $m$ 或 $n$ 同奇偶，并且
\[
\left|J(P_m,Q_n)\right|\leq\min\{m,n+1\}.
\]
又 $J(P_m,Q_n)$ 与 $n+1$ 奇偶性不同，故
\[
\left|J(P_m,Q_n)\right|\leq\min\{m,n\}.
\]
当 $Q_n(1,0)\ne0$ 时，同理可证。
\end{proof}

\begin{example}{}{homogeneous-index}
求平面微分方程
\[
\begin{cases}
\dfrac{\mathrm{d}x}{\mathrm{d}t}
=(x^2+xy+y^2)^k\displaystyle\prod_{i=1}^{s}(x-2iy)
=P_{2k+s}(x,y),\\[6pt]
\dfrac{\mathrm{d}y}{\mathrm{d}t}
=\pm(x^2-xy+y^2)^l\displaystyle\prod_{i=1}^{s}
\bigl[x-(2i-1)y\bigr]
=Q_{2l+s}(x,y)
\end{cases}
\]
的奇点 $x=y=0$ 的指数。
\end{example}

\begin{solution}
由于 $P_{2k+s}(1,0)=1\ne0$，由定理
\ref{thm:homogeneous-index}，
\[
J(P_{2k+s},Q_{2l+s})
=N\bigl(P_{2k+s}(t,1),Q_{2l+s}(t,1)\bigr).
\]
其中
\[
\begin{aligned}
P_{2k+s}(t,1)
&=(t-2)(t-4)\cdots(t-2s)(t^2+t+1)^k,\\
Q_{2l+s}(t,1)
&=\pm(t-1)(t-3)\cdots\bigl(t-(2s-1)\bigr)(t^2-t+1)^l.
\end{aligned}
\]
函数 $Q_{2l+s}(t,1)/P_{2k+s}(t,1)$ 的无穷间断点为
$t=2,4,\ldots,2s$，并且
\[
\varepsilon_{2i}^{-}=(-1)^{1+2(s-i)},\qquad
\varepsilon_{2i}^{+}=(-1)^{2(s-i)},\qquad i=1,\ldots,s.
\]
因此
\[
J(P_{2k+s},Q_{2l+s})
=\dfrac{1}{2}\sum_{i=1}^{s}
\left(\varepsilon_{2i}^{+}-\varepsilon_{2i}^{-}\right)=s.
\]
\end{solution}

由此可见，对于任意给定的自然数 $m,n$，当 $m+n$ 为偶数时，任取与
$m$ 或 $n$ 同奇偶且满足 $|J|\leq\min\{m,n\}$ 的整数 $J$，总可构造
齐次多项式 $P_m,Q_n$，使式~(2.3) 的指数为 $J$。

% ============================================================
%  §5  临界奇点指数的有理计算
% ============================================================
+\section{临界奇点指数的有理计算}

如果方程组~(2.1) 在奇点 $M_0=(x_0,y_0)$ 上的主方程~(2.3) 不具有孤立奇点
$x=y=0$，那么方程组~(2.1) 的奇点 $M_0$ 叫临界奇点。

\begin{example}{}{critical-point-linear}
给定微分方程组
\[
\begin{cases}
\dfrac{\mathrm{d}x}{\mathrm{d}t}=ax+by+X_2(x,y),\\
\dfrac{\mathrm{d}y}{\mathrm{d}t}=cx+dy+Y_2(x,y),
\end{cases} \tag{5.1}
\]
其中 $X_2,Y_2$ 包含二次以上的项。当 $ad-bc=0$ 时，式~(5.1) 在奇点
$M_0=(0,0)$ 上的主方程不具有孤立奇点 $M_0$，这时奇点 $M_0=(0,0)$
叫做式~(5.1) 的临界奇点。
\end{example}

\begin{example}{}{critical-point-example}
给定微分方程
\[
\begin{cases}
\dfrac{\mathrm{d}x}{\mathrm{d}t}=(x-1)(y-1),\\
\dfrac{\mathrm{d}y}{\mathrm{d}t}=(y-1)^2-(x-1)^4.
\end{cases} \tag{5.2}
\]
式~(5.2) 的奇点 $M_0=(1,1)$ 上的主方程为
\[
\dfrac{\mathrm{d}x}{\mathrm{d}t}=xy,\qquad
\dfrac{\mathrm{d}y}{\mathrm{d}t}=y^2,
\]
不具有孤立奇点 $x=y=0$，这时奇点 $M_0=(1,1)$ 叫做式~(5.2) 的临界奇点。
\end{example}

对于临界奇点，定理~4.1 不再成立。如何计算临界奇点的指数，过去还没有人讨论过。
下面也是高维新提供的关于临界奇点指数的有理计算方法。

给定 $x,y$ 的 $i$ 次齐次多项式 $P_i(x,y)$ $(i=m,m+1,\ldots,N_1)$ 与
$Q_i(x,y)$ $(i=n,n+1,\ldots,N_2)$。如果平面微分方程组
\[
\begin{cases}
\dfrac{\mathrm{d}x}{\mathrm{d}t}=P(x,y)=\displaystyle\sum_{i=m}^{N_1}P_i(x,y),\\[4pt]
\dfrac{\mathrm{d}y}{\mathrm{d}t}=Q(x,y)=\displaystyle\sum_{i=n}^{N_2}Q_i(x,y)
\end{cases} \tag{5.3}
\]
在圆周 $x^2+y^2=\rho^2$ $(\rho>0)$ 上无奇点，而在圆 $x^2+y^2<\rho^2$
内的所有奇点都是孤立的，此时 $P(0,\rho)$ 与 $Q(0,\rho)$ 不全为零，
不失一般性可假定 $P(0,\rho)\ne0$。

\begin{theorem}{}{critical-index-rational}
方程组~(5.3) 在圆 $x^2+y^2<\rho^2$ 内的所有奇点的指数和为
\[
-\dfrac12N\left[
\sum_{i=m}^{N_1}\rho^i(t^2+1)^{N_1-i}P_i(2t,t^2-1),
\sum_{i=n}^{N_2}\rho^i(t^2+1)^{N_2-i}Q_i(2t,t^2-1)
\right].
\]
\end{theorem}

\begin{proof}
事实上，上述所有奇点的指数和为 $I$，它等于由式~(5.3) 所确定的向量场
沿圆周 $x^2+y^2=\rho^2$ 的旋转数。在圆周上令
\[
x=\rho\sin\varphi,\qquad y=-\rho\cos\varphi,\qquad -\pi<\varphi<\pi,
\]
并作代换 $t=\tan\dfrac{\varphi}{2}$，于是
\[
\cos\varphi=\dfrac{1-t^2}{t^2+1},\qquad
\sin\varphi=\dfrac{2t}{t^2+1}.
\]
由计算旋转数的公式~(1.2)，得
\[
\begin{aligned}
I
&=-\dfrac{1}{2\pi}\int_{-\infty}^{+\infty}
\mathrm{d}\tan^{-1}
\dfrac{
Q\left(\dfrac{2t\rho}{t^2+1},\dfrac{t^2-1}{t^2+1}\rho\right)}
{
P\left(\dfrac{2t\rho}{t^2+1},\dfrac{t^2-1}{t^2+1}\rho\right)
}\\
&=-\dfrac{1}{2\pi}\int_{-\infty}^{+\infty}
\mathrm{d}\tan^{-1}
\dfrac{
\sum_{i=n}^{N_2}\rho^i(t^2+1)^{N_2-i}Q_i(2t,t^2-1)}
{
\sum_{i=m}^{N_1}\rho^i(t^2+1)^{N_1-i}P_i(2t,t^2-1)}.
\end{aligned}
\]
由于 $P(0,\rho)\ne0$，分母中不含 $t$ 而只含 $\rho$ 的项非零。将分子、
分母同乘适当的 $(t^2+1)$ 的幂后，分母与分子的次数之差不为正数。
由第3节性质~2 中公式~(3.2)，即得定理结论。
\end{proof}

\begin{theorem}{}{critical-index-homogeneous}
如果方程
\[
\begin{cases}
\dfrac{\mathrm{d}x}{\mathrm{d}t}=P_{N_1}(x,y),\\
\dfrac{\mathrm{d}y}{\mathrm{d}t}=Q_{N_2}(x,y)
\end{cases} \tag{5.4}
\]
与方程~(5.3) 的所有奇点都孤立，那么式~(5.3) 在全平面上的所有奇点的
指数和等于式~(5.4) 的唯一奇点 $x=y=0$ 的指数
$J(P_{N_1}(x,y),Q_{N_2}(x,y))$。
\end{theorem}

\begin{proof}
既然 $x=y=0$ 是式~(5.4) 的孤立奇点，式~(5.4) 不可能有其他奇点。
若式~(5.4) 还有奇点 $(x_0,y_0)\ne(0,0)$，则对任何实数 $k$ 有
\[
P_{N_1}(kx_0,ky_0)^2+Q_{N_2}(kx_0,ky_0)^2=0,
\]
这样奇点 $x=y=0$ 就不是孤立的。故式~(5.4) 只有唯一奇点。

式~(5.3) 在有界区域内只有有限个奇点，设为 $M_i$ $(i=1,2,\ldots,s)$。
构造从 $(P_{N_1},Q_{N_2})$ 到
$\left(\sum_{i=m}^{N_1}P_i,\sum_{i=n}^{N_2}Q_i\right)$ 的连续形变
\[
\left(
P_{N_1}+\lambda\sum_{i=m}^{N_1-1}P_i,\quad
Q_{N_2}+\lambda\sum_{i=n}^{N_2-1}Q_i
\right),\qquad 0\leq\lambda\leq1.
\]
取充分大的 $k$，使式~(5.3) 的奇点都在圆 $x^2+y^2<k^2$ 内，并使得当
$x^2+y^2=k^2$、$0\leq\lambda\leq1$ 时，上式两分量不同时为零。由同伦不变性，
\[
\gamma\bigl((P_{N_1},Q_{N_2}),x^2+y^2=k^2\bigr)
=\gamma\bigl((P,Q),x^2+y^2=k^2\bigr).
\]
再由第2节性质~1，左端为 $J(P_{N_1},Q_{N_2})$，右端为
$\sum_{i=1}^sJ_{M_i}(P,Q)$，定理得证。
\end{proof}

通常临界奇点的指数可借助定理~5.1 或定理~5.2 来计算。

\begin{example}{}{critical-index-zero}
求微分方程
\[
\dfrac{\mathrm{d}x}{\mathrm{d}t}=xy=P(x,y),\qquad
\dfrac{\mathrm{d}y}{\mathrm{d}t}=y^2+x^4=Q(x,y) \tag{5.5}
\]
的奇点 $M_0(0,0)$ 的指数。
\end{example}

\begin{solution}
上方程在奇点 $M_0(0,0)$ 上的主方程为
\[
\dfrac{\mathrm{d}x}{\mathrm{d}t}=xy,\qquad
\dfrac{\mathrm{d}y}{\mathrm{d}t}=y^2.
\]
它的奇点 $x=y=0$ 不是孤立的，故 $M_0$ 是式~(5.5) 的临界奇点。这时
计算 $M_0$ 的指数不能应用定理~2.1 或定理~4.1。

由于式~(5.5) 只有孤立奇点 $M_0$，故可应用定理~5.1。因 $Q(0,1)\ne0$，
\[
\begin{aligned}
J_{M_0}
&=\dfrac12N\left(
\sum_{i=2}^{4}(t^2+1)^{4-i}Q_i(2t,t^2-1),P_2(2t,t^2-1)
\right)\\
&=\dfrac12N(t^8+14t^4+1,2t^3-2t)=0.
\end{aligned}
\]
请读者画出方程~(5.5) 在 $M_0(0,0)$ 附近的相图来验证结论 $J_{M_0}=0$。
\end{solution}

\begin{example}{}{critical-index-two}
求微分方程
\[
\dfrac{\mathrm{d}x}{\mathrm{d}t}=xy=P(x,y),\qquad
\dfrac{\mathrm{d}y}{\mathrm{d}t}=y^2-x^4=Q(x,y) \tag{5.6}
\]
的奇点 $M_0(0,0)$ 的指数。
\end{example}

\begin{solution}
$M_0$ 是式~(5.6) 的临界奇点，也是式~(5.6) 的孤立奇点，故可应用定理~5.1
来计算 $M_0$ 的指数。因 $Q(0,1)\ne0$，由定理~5.1，
\[
\begin{aligned}
J_{M_0}
&=\dfrac12N\left(
\sum_{i=2}^{4}(t^2+1)^{4-i}Q_i(2t,t^2-1),P_2(2t,t^2-1)
\right)\\
&=\dfrac12N\left((t^2+1)^2(t^2-1)^2-(2t)^4,2t(t^2-1)\right)\\
&=\dfrac12N(t^8-18t^4+1,2t^3-2t)=2.
\end{aligned}
\]
请读者以第二章第3节中图~2.20 来验证 $J_{M_0}=2$。
\end{solution}

% ============================================================
%  §6  Bendixson公式
% ============================================================
\section{Bendixson公式}

给定微分方程
\[
\dfrac{\mathrm{d}x}{\mathrm{d}t}=X(x,y),\qquad
\dfrac{\mathrm{d}y}{\mathrm{d}t}=Y(x,y). \tag{6.1}
\]
设 $X(x,y),Y(x,y)$ 在原点 $O$ 的邻域中是 $x,y$ 的解析函数，原点 $O$
是式~(6.1) 的孤立奇点，$h,c,p$ 分别是原点 $O$ 的充分小邻域中双曲扇形、
椭圆扇形和抛物扇形的个数，又设奇点 $O$ 的指数是 $J$，则有
Bendixson 公式：
\[
J=1+\dfrac{c-h}{2}. \tag{6.2}
\]

\begin{proof}
证明 Bendixson 公式。

1. 设 $\{O\}$ 是中心型奇点，即 $\{O\}$ 的任意小邻域内部都有闭轨。
这时显然有 $J=1$，$c=h=0$，公式~(6.2) 自然成立。

2. 设 $\{O\}$ 是结点或焦点，即 $\{O\}$ 的充分小邻域内的每条轨线当
$t\to+\infty$（或 $-\infty$）时都进入奇点 $O$。对充分小 $\delta>0$，
可作在 $\partial S_\delta(O)$ 上与式~(6.1) 同伦的向量场，使得
$\partial S_\delta(O)$ 无切。于是其旋转数为 $1$，且 $c=h=0$。

3. 设 $O$ 为非中心型和非结点、焦点型奇点。作 $S_\delta(O)$ 满足第二章
第6节中的定理~6.1 和~6.3。于是由第二章第6节中的定理~6.2，
$S_\delta(O)$ 中必有轨线进入奇点 $O$。由于 $O$ 不是结点和焦点，
$S_\delta(O)$ 中至少有一个双曲扇形或椭圆扇形，又由于 $X(x,y),Y(x,y)$
解析，它们的个数有限。取 $\delta$ 充分小，使 $\partial S_\delta(O)$
与所有椭圆扇形相交。令 $L_1,L_2,\ldots,L_n$ 是不同扇形的边界，且逆时针
方向排列。设 $L_r$ 与圆周 $\partial S_\delta(O)$ 的最后一个交点为 $P_r$。
图~3.6 是 $h=2,c=2,p=2,n=6$ 时的示意图。
\end{proof}

\chThreeFigurePlaceholder{3.6}

对抛物扇形 $\widehat{OP_kP_{k+1}}$，将圆弧
$\widehat{P_kP_{k+1}}$ 叫做抛物弧段；对双曲扇形
$\widehat{OP_iP_{i+1}}$，取轨线段 $\widehat{P'_iT_iP'_{i+1}}$
叫做双曲弧段；对椭圆扇形 $\widehat{OP_jP_{j+1}}$，取轨线段
$\widehat{P'_jT_jP'_{j+1}}$ 叫做椭圆弧段。将这些弧段与圆弧段顺次连接
成一条光滑闭曲线 $l$，并在 $l$ 上把向量场连续形变为与相应弧段相切或法向一致
的向量场 $A_1(x,y)$。由同伦不变性，$A_1$ 与原向量场在 $l$ 上有相同的旋转数。

对双曲扇形，若两边与 $O$ 的连线夹角为 $\alpha_i$，则 $A_1$ 沿相应弧段的
旋转角为 $\alpha_i-\pi$；对椭圆扇形，夹角为 $\beta_i$，旋转角为
$\beta_i+\pi$；对抛物扇形，夹角为 $\gamma_k$，旋转角为 $\gamma_k$。
又
\[
\sum_{i=1}^{h}\alpha_i+\sum_{j=1}^{c}\beta_j+
\sum_{k=1}^{p}\gamma_k=2\pi. \tag{6.3}
\]
故
\[
\gamma(A_1,l)
=\dfrac{1}{2\pi}\left[
\sum_{i=1}^{h}(\alpha_i-\pi)+
\sum_{j=1}^{c}(\beta_j+\pi)+
\sum_{k=1}^{p}\gamma_k\right]
=1+\dfrac{c-h}{2}.
\]
另一方面，向量场 $A_0(x,y)=(X(x,y),Y(x,y))$ 沿闭曲线 $l$ 的旋转数就是
奇点 $O$ 的指数。因 $A_0$ 与 $A_1$ 在 $l$ 上同伦，
\[
J(A_0)=\gamma(A_0,l)=\gamma(A_1,l)
=1+\dfrac{c-h}{2}.
\]
Bendixson 公式证毕。

\noindent\textbf{附注}\quad
Bendixson 公式对于连续向量场也是成立的。只是要将 $c$ 换成与边界
$\partial S_\delta(O)$ 至少有两个公共点（不是切点）的椭圆扇形的个数，
要将 $h$ 换成双曲扇形及双曲椭圆扇形个数的和。

% ============================================================
%  习题和参考文献
% ============================================================
\section*{习题和参考文献}
\subsection*{习题}

\begin{enumerate}
\item 试证经坐标轴的移动和旋转，连续向量场的旋转数不变。
\item 试证经反射变换 $(x,y)\mapsto(-x,y)$ 或 $(x,y)\mapsto(x,-y)$，
连续向量场的旋转数相差一个符号。
\item 设连续向量场 $A(x,y)=(X(x,y),Y(x,y))$ 在区域 $D$ 的边界
$\partial D$ 上非退化，试证
\[
\gamma((X,Y),\partial D)=-\gamma((-X,Y),\partial D).
\]
\item 考虑微分方程
\[
\dfrac{\mathrm{d}x}{\mathrm{d}t}=ax+by,\qquad
\dfrac{\mathrm{d}y}{\mathrm{d}t}=cx+dy.
\]
试证当行列式
\[
\left|\begin{matrix}a&b\\c&d\end{matrix}\right|\ne0
\]
时，有
\[
J(ax+by,cx+dy)
=\operatorname{sgn}\left|\begin{matrix}a&b\\c&d\end{matrix}\right|J(x,y)
=\operatorname{sgn}\left|\begin{matrix}a&b\\c&d\end{matrix}\right|.
\]
由此推得，当 $O(0,0)$ 是鞍点时，奇点指数为 $-1$；当 $O(0,0)$ 是结点、
中心、焦点时，奇点指数为 $1$。
\item 给定平面微分方程
\[
\dfrac{\mathrm{d}x}{\mathrm{d}t}=ax+by+h(x,y),\qquad
\dfrac{\mathrm{d}y}{\mathrm{d}t}=cx+dy+g(x,y),
\]
其中 $ad\ne bc$，又 $h(x,y),g(x,y)$ 在 $x=y=0$ 附近有连续偏导数，并设
\[
\lim_{x^2+y^2\to0}\dfrac{h(x,y)}{\sqrt{x^2+y^2}}
=\lim_{x^2+y^2\to0}\dfrac{g(x,y)}{\sqrt{x^2+y^2}}=0.
\]
证明此方程的奇点 $x=y=0$ 的指数为 $\operatorname{sgn}(ad-bc)$。由此推得，
当
\[
\left|\begin{matrix}a&b\\c&d\end{matrix}\right|=ad-bc\ne0
\]
时，非线性方程的奇点指数等于相应的线性方程的奇点指数。
\item $u(t)=t^4+3t^3+2t^2-2t+1$，$v(t)=t^2-t+1$，求 $N(u,v)$ 及
$N(v,u)$，并验证 Cauchy 指标的性质~1、3、4。
\item 叫多项式 $f(t)$ 为正定（或负定）的，如果存在正数 $\varepsilon>0$，
使得 $f(t)\geq\varepsilon$（或 $f(t)\leq-\varepsilon$）。正定的或负定的多项式
叫做定号多项式。试证若 $g_1(t),g_2(t)$ 为定号多项式，则
\[
N(g_1(t)u(t),g_2(t)v(t))
=\begin{cases}
N(u,v),&\text{当 }g_1,g_2\text{ 同为正定或负定时},\\
-N(u,v),&\text{当 }g_1,g_2\text{ 一为正定另一为负定时}.
\end{cases}
\]
\item 对于两个无公共间断点的分式函数来说，它们的 Cauchy 指标的和等于它们和的
Cauchy 指标。
\item 求平面微分方程
\[
\begin{cases}
\dfrac{\mathrm{d}x}{\mathrm{d}t}=x^3+x^2y-xy^2+y^3+x^4,\\
\dfrac{\mathrm{d}y}{\mathrm{d}t}=x^5+x^3y^2-x^2y^3+y^5+y^6
\end{cases}
\]
的奇点 $O(0,0)$ 的指数。
\item 求平面微分方程
\[
\dfrac{\mathrm{d}x}{\mathrm{d}t}=xy-3x^3,\qquad
\dfrac{\mathrm{d}y}{\mathrm{d}t}=y^2-6x^2y+x^4
\]
的奇点 $O(0,0)$ 的指数。
\item 求平面微分方程
\[
\dfrac{\mathrm{d}x}{\mathrm{d}t}=2x^3-y,\qquad
\dfrac{\mathrm{d}y}{\mathrm{d}t}=y^2+x^s\quad(s\geq3)
\]
的奇点 $O(0,0)$ 的指数。
\item 试证在结点附近可作包围结点的无切单闭光滑曲线。
\item 对例题~5.3、5.4 和习题~10、11 验证 Bendixson 公式。
\end{enumerate}

\subsection*{参考文献}

\begin{enumerate}
\item 高维新，平面奇点的指数，微分方程论文集，北京大学数学力学系，1963 年，189--198。
\item 李正元、钱敏，向量场的旋转度理论及其应用，北京大学出版社，1982。
\item 刘品善，关于微分方程高次奇点的指数，北京大学学报，4(1957)，395--401。
\item Коренцкий, П.\ Т., Общий метод вычисления индекса особой точки
дифференциального уравнения, \textit{Научн. Зап. Одес. Политехн. Ин-та},
23(1960), 56--62.
\end{enumerate}

\endinput
